% Pedro. S1--S5. Contexto, exchange, order book y procedimiento actual de
% recuperación en Binance.
\section{Contexto}
\begin{frame}
  \vspace{5mm}
  {\LARGE\color{colAcento}\bfseries RelayBook\par}
  \vspace{3mm}
  {\large\color{colAcento}Mejora de la disponibilidad del market data de Binance\\ante \emph{reconnection storms} con \emph{state-based resume}\par}
  \vspace{8mm}
  Nicolas Castaneda y Pedro Mori\par
  Arquitectura de Software --- UTEC 2026-2\par
  Noviembre de 2026
  \note{Pedro. 0:00--0:15 (15 s). Buenos días. La presentación expone RelayBook, una
  arquitectura para que los clientes de Binance recuperen sus datos de mercado
  sin saturar al servidor central cuando miles se desconectan a la vez.}
\end{frame}

\begin{frame}{Binance como \emph{exchange}}
  \centering% S2. Qué es un exchange: órdenes que entran, libro y matching dentro,
% market data que sale hacia millones de clientes.
\begin{tikzpicture}[x=1cm,y=1cm]
  \path[use as bounding box] (0,0.1) rectangle (13.5,4.75);
  % Personas que mandan órdenes
  \node[anchor=west,font=\small\bfseries] at (0,4.5) {Personas y bots mandan órdenes};
  \node[box,font=\scriptsize,text width=2.9cm,text=colVenta,draw=colVenta] (o1) at (1.7,3.6) {Ana: \textbf{vendo} 0.5 BTC\\a 67\,002 USD};
  \node[box,font=\scriptsize,text width=2.9cm,text=colCompra,draw=colCompra] (o2) at (1.7,2.3) {Luis: \textbf{compro} 1.2 BTC\\a 67\,000 USD};
  \node[box,font=\scriptsize,text width=2.9cm,text=colCompra,draw=colCompra] (o3) at (1.7,1.0) {Bot: \textbf{compro} 1 BTC\\ya, al mejor precio};
  % Exchange
  \draw[colAcento,line width=1pt,rounded corners=4pt,fill=colAcento!4] (4.4,0.4) rectangle (8.9,4.25);
  \node[font=\small\bfseries,text=colAcento] at (6.65,3.95) {Binance (exchange)};
  \node[box,font=\scriptsize,text width=3.6cm] (me) at (6.65,3.0) {\textbf{Matching engine}\\cruza compras con ventas};
  \node[acento,font=\scriptsize,text width=3.6cm] (ob) at (6.65,1.4) {\textbf{Order book}\\las órdenes que todavía\\no encuentran pareja};
  \draw[flow,{Stealth[length=2.3mm]}-{Stealth[length=2.3mm]}] (me) -- (ob);
  \foreach \o/\y in {o1/3.6,o2/2.3,o3/1.0} \draw[flow] (\o.east) -- (4.4,\y);
  % Market data
  \node[anchor=east,font=\small\bfseries] at (13.5,4.5) {Market data};
  \node[box,font=\scriptsize,text width=2.7cm] (c1) at (11.85,3.6) {Apps de trading};
  \node[box,font=\scriptsize,text width=2.7cm] (c2) at (11.85,2.3) {Bots y market makers};
  \node[box,font=\scriptsize,text width=2.7cm] (c3) at (11.85,1.0) {Proveedores de datos};
  \foreach \c/\y in {c1/3.6,c2/2.3,c3/1.0} \draw[flow,colAcento] (8.9,\y) -- (\c.west);
  \node[font=\scriptsize,text=colAcento,align=center] at (11.85,0.35) {millones de clientes};
\end{tikzpicture}

  \takeaway{El precio surge del cruce entre órdenes de compra y de venta.}
  \note{Pedro. 0:15--0:50 (35 s). Binance es el exchange de criptomonedas más grande
  del mundo. Un exchange es un mercado donde las personas compran y venden
  entre sí a través de la plataforma. En la figura, Ana ofrece vender medio
  bitcoin a 67 002 dólares y Luis ofrece comprar 1.2 a 67 000. Las órdenes sin
  contraparte esperan en el order book, y el matching engine cruza las que
  coinciden. A la derecha, Binance publica lo que ocurre en el mercado, el
  market data, para millones de apps, bots y proveedores.}
\end{frame}

\begin{frame}{El \emph{order book}}
  \centering% S3. Order book de BTC/USDT: dos lados, niveles de precio y qué es un delta.
\begin{tikzpicture}[x=1cm,y=1cm]
  \path[use as bounding box] (0,-0.1) rectangle (13.5,4.75);
  % ---------- Libro ----------
  \node[anchor=west,font=\scriptsize\bfseries,text=colVenta] at (0.6,4.5) {VENTAS (\emph{asks})};
  \node[venta] (a2) at (2.3,3.85) {67\,002 USD \quad 0.5 BTC};
  \node[venta] (a1) at (2.3,3.25) {67\,001 USD \quad 2.0 BTC};
  \draw[colNeutro,dashed,thick] (0.4,2.75) -- (4.2,2.75);
  \node[compra] (b1) at (2.3,2.25) {67\,000 USD \quad 1.2 BTC};
  \node[compra] (b2) at (2.3,1.65) {66\,999 USD \quad 0.8 BTC};
  \node[anchor=west,font=\scriptsize\bfseries,text=colCompra] at (0.6,1.05) {COMPRAS (\emph{bids})};
  % Precio de mercado
  \node[font=\scriptsize,text=colAcento,align=left,anchor=west] (pm) at (4.45,2.75) {precio de mercado\\entre 67\,000 y 67\,001};
  % Nivel = suma de órdenes
  \node[font=\scriptsize,align=left,anchor=north west,text width=6.5cm] at (0.1,0.7) {Cada fila es un \textbf{nivel de precio} y suma todas las órdenes a ese precio, por ejemplo 2.0 = 0.5 + 1.5 de dos vendedores.};
  % ---------- Qué cambia el libro ----------
  \draw[colNeutro!40] (7.0,-0.05) -- (7.0,4.7);
  \node[anchor=west,font=\small\bfseries] at (7.3,4.5) {Qué cambia el libro};
  \node[box,font=\scriptsize,text width=5.4cm] (s1) at (10.35,3.75) {1. Llega una orden de compra de \textbf{1 BTC} al mejor precio};
  \node[box,font=\scriptsize,text width=5.4cm] (s2) at (10.35,2.75) {2. La orden se cruza con el vendedor más barato en \textbf{67\,001}};
  \node[box,font=\scriptsize,text width=5.4cm] (s3) at (10.35,1.75) {3. El nivel 67\,001 pasa de 2.0 a \textbf{1.0 BTC}};
  \node[acento,font=\scriptsize,text width=5.4cm] (s4) at (10.35,0.6) {Ese cambio es un \textbf{delta}, que trae el valor nuevo\\\emph{``67\,001 ahora vale 1.0''}};
  \draw[flow] (s1) -- (s2);
  \draw[flow] (s2) -- (s3);
  \draw[flow] (s3) -- (s4);
\end{tikzpicture}

  \takeaway{Cada orden nueva, cancelada o ejecutada cambia un nivel del libro.}
  \note{Pedro. 0:50--2:05 (75 s). En el order book las ventas, o asks, aparecen
  arriba en ámbar y las compras, o bids, abajo en morado. El vendedor más
  barato pide 67 001 y el comprador que más paga ofrece 67 000, por lo que el
  precio de mercado queda entre ambos valores. Cada fila es un nivel de precio
  y suma todas las órdenes a ese precio, como el 2.0 de 67 001, que
  corresponde a dos vendedores con 0.5 y 1.5. A la derecha, un bot compra 1
  bitcoin al mejor precio. La orden se cruza con el vendedor de 67 001 y ese
  nivel baja de 2.0 a 1.0. Ese cambio se llama delta y contiene el valor nuevo
  del nivel, 1.0, en lugar de una resta. La propuesta se apoya en esta
  propiedad. Binance publica el libro nivel por nivel con su cantidad total, y
  con eso cada cliente arma una copia local.}
\end{frame}

\begin{frame}{Copia local con \emph{snapshot} y deltas}
  \centering% S5. Cómo el cliente arma su copia: snapshot + deltas numerados, y qué pasa
% con la copia si falta un mensaje. Mismo libro de la S3 y la S4.
\begin{tikzpicture}[x=1cm,y=1cm]
  \path[use as bounding box] (0,-0.3) rectangle (13.5,5.45);
  \def\pv#1{\textcolor{colVenta}{#1}}
  \def\pc#1{\textcolor{colCompra}{#1}}
  \tikzset{
    libro/.style={box,font=\scriptsize,text width=2.0cm,inner sep=3pt},
    delta/.style={box,font=\scriptsize,text width=1.25cm,inner sep=3pt},
    def/.style={font=\scriptsize,align=left,anchor=north west,text width=4.2cm},
    fila/.style={font=\scriptsize\bfseries,anchor=west},
  }
  % ---------- Definiciones ----------
  \node[def] at (0.0,5.45) {El \textbf{snapshot} es la foto completa del libro en un momento, con 2\,000 niveles y 41\,KB en la grabación.};
  \node[def] at (4.6,5.45) {Un \textbf{delta} contiene solo los niveles que cambiaron, con su valor nuevo.};
  \node[def] at (9.2,5.45) {Cada delta lleva un \textbf{número de secuencia} en orden, como \#1001 y \#1002.};
  % ---------- Fila 1: llegan todos ----------
  \node[fila,text=colReq] at (0.0,4.3) {Si llegan todos los mensajes};
  \node[libro,acento] (s) at (1.2,3.25) {\textbf{Snapshot \#1000}\\[1pt]
    \begin{tabular}{@{}l@{\ \ }r@{}}\pv{67\,002}&0.5\\\pv{67\,001}&2.0\\\pc{67\,000}&1.2\\\pc{66\,999}&0.8\end{tabular}};
  \node[delta] (d1) at (3.4,3.25) {\textbf{\#1001}\\67\,001\\$\to$ 1.0};
  \node[delta] (d2) at (5.3,3.25) {\textbf{\#1002}\\67\,000\\$\to$ 1.5};
  \node[delta] (d3) at (7.2,3.25) {\textbf{\#1003}\\66\,999\\$\to$ 0\\(se borra)};
  \node[libro,bueno] (r) at (9.6,3.25) {\textbf{Copia del cliente}\\[1pt]
    \begin{tabular}{@{}l@{\ \ }r@{}}\pv{67\,002}&0.5\\\pv{67\,001}&\textbf{1.0}\\\pc{67\,000}&\textbf{1.5}\\\pc{66\,999}&---\end{tabular}};
  \node[bueno,font=\scriptsize,text width=1.9cm] (ok) at (12.25,3.25) {igual al\\libro real \okmark};
  \foreach \a/\b in {s/d1,d1/d2,d2/d3,d3/r} \draw[flow] (\a) -- (\b);
  \draw[flow,colReq] (r) -- (ok);
  % ---------- Fila 2: falta el #1002 ----------
  \node[fila,text=colWake] at (0.0,2.0) {Si se pierde el \#1002};
  \node[libro,acento] (s') at (1.2,0.95) {\textbf{Snapshot \#1000}\\[1pt]
    \begin{tabular}{@{}l@{\ \ }r@{}}\pv{67\,002}&0.5\\\pv{67\,001}&2.0\\\pc{67\,000}&1.2\\\pc{66\,999}&0.8\end{tabular}};
  \node[delta] (e1) at (3.4,0.95) {\textbf{\#1001}\\67\,001\\$\to$ 1.0};
  \node[delta,draw=colWake,text=colWake,dashed,fill=white] (e2) at (5.3,0.95) {\textbf{\#1002}\\no llegó};
  \node[delta] (e3) at (7.2,0.95) {\textbf{\#1003}\\66\,999\\$\to$ 0\\(se borra)};
  \node[libro,malo] (w) at (9.6,0.95) {\textbf{Copia del cliente}\\[1pt]
    \begin{tabular}{@{}l@{\ \ }r@{}}\pv{67\,002}&0.5\\\pv{67\,001}&1.0\\\pc{67\,000}&\textbf{1.2} \xmark\\\pc{66\,999}&---\end{tabular}};
  \node[malo,font=\scriptsize,text width=1.9cm] (bad) at (12.25,0.95) {dice 1.2 y el real es 1.5, \textbf{no confiable}};
  \draw[flow] (s') -- (e1);
  \draw[flow,colWake,dashed] (e1) -- (e2);
  \draw[flow,colWake,dashed] (e2) -- (e3);
  \draw[flow] (e3) -- (w);
  \draw[flow,colWake] (w) -- (bad);
  \node[font=\scriptsize,text=colWake,anchor=north] at (5.3,0.05) {el \#1003 llega después del \#1001, así que el cliente detecta la pérdida};
\end{tikzpicture}

  \note{Pedro. 2:05--2:55 (50 s). El cliente arma esa copia con dos tipos de
  mensaje. El snapshot es la foto completa del libro, que en la grabación
  tenía 2 000 niveles y 41 KB. Los deltas contienen solo los niveles que
  cambiaron y llevan un número de secuencia. En la fila superior llegan el
  1001, el 1002 y el 1003, y la copia queda igual al libro real. En la fila
  inferior se pierde el 1002, por lo que la copia muestra 1.2 en el nivel 67
  000 cuando el valor real es 1.5. El cliente detecta la pérdida porque
  después del 1001 recibe el 1003, y desde ese momento su copia deja de ser
  confiable.}
\end{frame}

\begin{frame}{Recuperación en el contrato actual de Binance}
  \centering% S5. El contrato documentado de Binance para mantener un libro local.
\begin{tikzpicture}[x=1cm,y=1cm]
  \path[use as bounding box] (0,0.55) rectangle (13.5,4.45);
  \node[box,font=\scriptsize,text width=3.0cm] (n1) at (1.9,3.8) {\textbf{1.} Abrir el \textbf{WebSocket},\\conexión abierta que\\empuja los deltas};
  \node[box,font=\scriptsize,text width=3.0cm] (n2) at (5.6,3.8) {\textbf{2.} Pedir el snapshot por \textbf{REST}\\al servidor central};
  \node[box,font=\scriptsize,text width=3.0cm] (n3) at (9.3,3.8) {\textbf{3.} Aplicar los deltas\\sobre el snapshot};
  \draw[flow] (n1) -- (n2);
  \draw[flow] (n2) -- (n3);
  \node[decision] (d) at (9.3,1.75) {¿salto en la\\secuencia\\o corte?};
  \draw[flow] (n3) -- (d);
  \draw[flow,colReq] (d.east) -- ++(1.0,0) |- (n3.east);
  \node[font=\scriptsize\bfseries,text=colReq,anchor=west] at (12.0,2.75) {no};
  \node[malo,font=\scriptsize,text width=3.0cm] (bad) at (5.6,1.75) {Tirar el libro y pedir\\otra vez el \textbf{snapshot completo}};
  \draw[flow,colWake] (d.west) -- node[above,font=\scriptsize\bfseries,text=colWake] {sí} (bad.east);
  \draw[flow,colWake] (bad.north) -- (n2.south);
  \node[malo,font=\scriptsize,text width=3.0cm] (lim) at (1.9,1.75) {Rate limit del REST\\exceso de peticiones $\to$ \textbf{429}\\insistencia $\to$ \textbf{418} (ban)};
  \draw[colWake,dashed,thick] (bad.west) -- (lim.east);
\end{tikzpicture}

  \takeaway{Sin reenvío de mensajes perdidos, un corte de 1\,s exige el libro completo.}
  \source{Documentación oficial de Binance, \emph{How to manage a local order book correctly}.}
  \note{Pedro. 2:55--3:35 (40 s). Binance documenta un procedimiento en tres pasos.
  Primero, el cliente abre un WebSocket, una conexión persistente que empuja
  los deltas. Luego pide el snapshot por REST al servidor central y aplica los
  deltas encima. Ante un salto en la secuencia o un corte, el cliente descarta
  su libro y vuelve a pedir el snapshot completo. Ese endpoint tiene rate
  limit, con error 429 ante exceso de peticiones y 418, un ban, si el cliente
  insiste. Como el stream no reenvía los mensajes perdidos, un corte de un
  segundo o de un minuto cuesta el libro entero.}
\end{frame}
